\documentclass[12pt, a4paper]{article}

\usepackage[utf8]{inputenc}
\usepackage[T1]{fontenc}
\usepackage[english]{babel}
\usepackage{amsmath, amssymb, amsthm}
\usepackage{geometry}
\usepackage{booktabs}
\usepackage{array}
\usepackage{xcolor}
\usepackage{hyperref}
\usepackage{algorithm}
\usepackage{algpseudocode}
\usepackage{parskip}
\usepackage{titlesec}
\usepackage{fancyhdr}
\usepackage{float}

\geometry{top=2.5cm, bottom=2.5cm, left=2.5cm, right=2.5cm}

\hypersetup{
    colorlinks=true,
    linkcolor=blue!60!black,
    urlcolor=blue!60!black,
    citecolor=blue!60!black
}

\titleformat{\section}{\large\bfseries}{\thesection}{1em}{}[\titlerule]
\titleformat{\subsection}{\normalsize\bfseries}{\thesubsection}{1em}{}

\pagestyle{fancy}
\fancyhf{}
\rhead{\small Interest Rate Modeling}
\lhead{\small Ecole Centrale Marseille}
\cfoot{\thepage}

\newcommand{\E}{\mathbb{E}}
\newcommand{\Q}{\mathbb{Q}}
\newcommand{\diff}{\mathrm{d}}

\begin{document}

% ===== TITLE PAGE =====
\begin{titlepage}
\centering
\vspace*{3cm}
{\LARGE\bfseries Hull-White Calibration and\\[0.3em]Multi-Measure Valuation\par}
\vspace{0.8cm}
{\large\itshape Interest Rate Modeling Project\par}
\vspace{3cm}
{\large Ecole Centrale Marseille\par}
\vspace{0.4cm}
{\normalsize Pauline DUMAS, Artouch VARDANYAN --- March 2026\par}
\vspace{3cm}
\begin{abstract}
\noindent
This report presents the calibration and application of the one-factor Hull-White model for pricing a European payer swaption. We estimate the historical parameters $(a, \sigma)$ from \texteuro STR data using two methods (AR(1) regression and exact MLE), fit the initial yield curve via Nelson-Siegel-Svensson (NSS), and price the swaption under three distinct probability measures ($\Q$, $\Q^T$, $\Q^A$) using Monte Carlo simulation. We then implement a recombining trinomial tree following Hull \& White (1993) and compare the results, analyze convergence, and study the sensitivity to the mean-reversion parameter. All methods converge to a consistent swaption price near $0.0990$ per unit notional.
\end{abstract}
\end{titlepage}

\tableofcontents

% ===== SECTION 1 =====
\section{Calibration Strategy}

\subsection{Model Overview}

The one-factor Hull-White model describes the short rate $r_t$ under the risk-neutral measure $\Q$ as:
\begin{equation}
\diff r_t = (\theta(t) - a\, r_t)\,\diff t + \sigma\, \diff W_t^{\Q}
\end{equation}
where $a > 0$ is the mean-reversion speed, $\sigma > 0$ the volatility, and $\theta(t)$ a deterministic function chosen so that the model exactly fits the initial yield curve. This is a time-inhomogeneous Ornstein-Uhlenbeck (OU) process.

The calibration has two separate steps. First, we estimate $a$ and $\sigma$ from historical short-rate data under the real-world measure $P$, assuming they are the same under both measures. Then, independently, we determine $\theta(t)$ from the current market yield curve under $\Q$. In other words, historical data is only used to pin down $a$ and $\sigma$; the pricing drift is entirely driven by $\theta(t)$.

\subsection{Historical Parameter Estimation}

\subsubsection{Data}

We use daily \texteuro STR (Euro Short-Term Rate) data from the ECB data portal, covering October 1, 2019 to February 19, 2026, giving $n = 1637$ observations with average time step $\Delta t \approx 0.0039$ years.

\subsubsection{AR(1) Regression Approach}

The exact discretization of the OU process over one step $\Delta t$ gives:
\begin{equation}
r_{t+\Delta t} = \underbrace{b(1-e^{-a\Delta t})}_{\alpha} + \underbrace{e^{-a\Delta t}}_{\beta}\,r_t + \varepsilon_t, \quad \varepsilon_t \sim \mathcal{N}\!\left(0,\, \frac{\sigma^2}{2a}(1-e^{-2a\Delta t})\right)
\end{equation}
This is an AR(1) regression in $(\hat{\alpha}, \hat{\beta})$. We fit it by OLS, then recover the structural parameters using the exact inversion formulas:
\begin{equation}
\hat{a} = -\frac{\ln\hat{\beta}}{\Delta t}, \quad \hat{b} = \frac{\hat{\alpha}}{1-\hat{\beta}}, \quad \hat{\sigma} = \hat{s}_\varepsilon\sqrt{\frac{2\hat{a}}{1-\hat{\beta}^2}}
\end{equation}

\begin{algorithm}[H]
\caption{AR(1) Estimation (exact discretization)}
\begin{algorithmic}[1]
\State Set $y_i = r_{t_i + \Delta t}$, $X_i = (1,\, r_{t_i})$
\State Estimate $(\hat{\alpha}, \hat{\beta}) = (X^\top X)^{-1} X^\top y$ by OLS
\State Compute residual std $\hat{s}_\varepsilon = \mathrm{std}(\hat{\varepsilon})$
\State $\hat{a}_{\text{AR}} \leftarrow -\ln(\hat{\beta}) / \Delta t$; \quad $\hat{b}_{\text{AR}} \leftarrow \hat{\alpha} / (1 - \hat{\beta})$; \quad $\hat{\sigma}_{\text{AR}} \leftarrow \hat{s}_\varepsilon\sqrt{2\hat{a}_{\text{AR}}/(1-\hat{\beta}^2)}$
\end{algorithmic}
\end{algorithm}

\subsubsection{Exact Maximum Likelihood Estimation (MLE)}

MLE uses the same Gaussian transition distribution, but instead of minimizing a sum of squares it maximizes the exact likelihood. Setting $\phi = e^{-a\Delta t}$, the closed-form estimators are:
\begin{align}
\hat{\phi} &= \frac{S_{xy} - S_x S_y / n}{S_{xx} - S_x^2 / n}, \quad
\hat{\alpha} = \frac{S_y - \hat{\phi}\, S_x}{n(1-\hat{\phi})}, \quad
\hat{a}_{\text{MLE}} = -\frac{\ln\hat{\phi}}{\Delta t}, \quad
\hat{b}_{\text{MLE}} = \frac{\hat{\alpha}}{1-\hat{\phi}}
\end{align}
The volatility is then recovered from the residual standard deviation:
\begin{equation}
\hat{\sigma}_{\text{MLE}} = \hat{s}_\varepsilon \sqrt{\frac{2\hat{a}}{1 - \hat{\phi}^2}}
\end{equation}

\begin{algorithm}[H]
\caption{Exact MLE for the OU Process}
\begin{algorithmic}[1]
\State Compute sufficient statistics: $S_x, S_y, S_{xx}, S_{xy}$ from the rate series
\State $\hat{\phi} \leftarrow (S_{xy} - S_x S_y/n) / (S_{xx} - S_x^2/n)$
\State $\hat{a}_{\text{MLE}} \leftarrow -\ln(\hat{\phi}) / \Delta t$; \quad $\hat{b}_{\text{MLE}} \leftarrow \hat{\alpha} / (1 - \hat{\phi})$
\State Compute residuals; $\hat{\sigma}_{\text{MLE}} \leftarrow \mathrm{std}(e)\cdot\sqrt{2\hat{a}/(1-\hat{\phi}^2)}$
\end{algorithmic}
\end{algorithm}

\subsubsection{Results and Comparison}

\begin{table}[H]
\centering
\begin{tabular}{lcc}
\toprule
\textbf{Parameter} & \textbf{AR(1)} & \textbf{MLE} \\
\midrule
$a$ (mean-reversion speed) & 0.13668 & 0.13672 \\
$\sigma$ (volatility)       & 0.00670 & 0.00670 \\
$b$ (long-run mean)         & 0.04093 & 0.04093 \\
\bottomrule
\end{tabular}
\caption{Estimated parameters from \texteuro STR data (Oct.\ 2019 -- Feb.\ 2026)}
\end{table}

Both methods give virtually identical estimates. We keep the MLE parameters for the rest of the project: $a = 0.1367$, $\sigma = 0.0067$, $b = 0.0409$.

\textbf{Discussion.} The two methods share the same exact OU discretization and the same inversion formulas, so numerically they agree completely — as confirmed by Table~1. The real difference is in how $(\hat{\alpha}, \hat{\beta})$ are obtained: AR(1) uses OLS, which minimizes a sum of squares without reference to the transition distribution; MLE directly maximizes the Gaussian likelihood and achieves the Cramér-Rao bound asymptotically. Both are fine here, but we prefer MLE on statistical grounds.

\subsection{Term Structure Fitting: Nelson-Siegel-Svensson}

\subsubsection{Market Data}

We use the ECB AAA-rated yield curve (\texttt{G\_N\_A}) as of March 5, 2026, on 16 maturities from 3 months to 30 years.

\subsubsection{The NSS Model}

The NSS model parameterizes the spot rate $z(0,T)$ as:
\begin{equation}
z^{\text{NSS}}(T) = \beta_0 + \beta_1 \frac{1-e^{-T/\tau_1}}{T/\tau_1} + \beta_2\!\left(\frac{1-e^{-T/\tau_1}}{T/\tau_1} - e^{-T/\tau_1}\right) + \beta_3\!\left(\frac{1-e^{-T/\tau_2}}{T/\tau_2} - e^{-T/\tau_2}\right)
\end{equation}
The instantaneous market forward rate $f^M(0,T)$ follows analytically:
\begin{equation}
f^M(0,T) = \beta_0 + \beta_1 e^{-T/\tau_1} + \beta_2 \frac{T}{\tau_1} e^{-T/\tau_1} + \beta_3 \frac{T}{\tau_2} e^{-T/\tau_2}
\end{equation}
To avoid local minima, calibration uses a grid search over $(\tau_1, \tau_2)$ starting points before running \texttt{least\_squares} (TRF method).

\subsubsection{Results}

\begin{table}[H]
\centering
\begin{tabular}{lcc}
\toprule
\textbf{Parameter} & \textbf{Value} & \textbf{Interpretation} \\
\midrule
$\beta_0$ & 0.012309 & Long-term level ($\approx 1.23\%$) \\
$\beta_1$ & 0.007979 & Slope \\
$\beta_2$ & $-$0.011718 & First curvature \\
$\beta_3$ & 0.076372 & Second curvature \\
$\tau_1$  & 4.751 years & \\
$\tau_2$  & 14.768 years & \\
\midrule
RMSE            & 0.033 bps & \\
Max abs.\ error & 0.060 bps & \\
\bottomrule
\end{tabular}
\caption{NSS calibrated parameters -- ECB AAA curve, 05/03/2026}
\end{table}

The fit is very good: the RMSE on spot rates is around $3.3 \times 10^{-5}$ in decimal, which is about $0.033$ basis points. The market ZCB price is $P^M(0,T) = e^{-T \cdot z^{\text{NSS}}(T)}$.

The parameters paint a consistent picture of the ECB curve on that date. The long-term level $\beta_0 \approx 1.23\%$ is relatively low, reflecting an environment where long-run rate expectations remain contained. The positive slope $\beta_1 > 0$ confirms that the curve is upward-sloping: short rates are below long rates, which is the normal shape. The dominant feature is $\beta_3 = 0.076$, by far the largest coefficient, which creates a pronounced hump around $\tau_2 \approx 15$ years — this is what gives the ECB AAA curve its characteristic belly in the intermediate maturities. The negative $\beta_2$ introduces a slight concavity at shorter maturities, consistent with the compression of spreads in the 1--5 year segment.

\subsubsection{Determination of $\theta(t)$}

For Hull-White to exactly fit the initial yield curve, $\theta(t)$ must be chosen so that the model-implied ZCB prices match the market at all maturities: $P^{\text{model}}(0,T) = P^M(0,T)$ for all $T$.

Under Hull-White, the ZCB price has the affine form $P(0,T) = \mathcal{A}(0,T)\,e^{-\mathcal{B}(0,T)\,r_0}$. Taking $-\frac{1}{T}\ln P(0,T)$ and differentiating with respect to $T$ gives the model-implied instantaneous forward rate $f^{\text{model}}(0,T)$. Matching it to the market forward rate $f^M(0,T)$ leads to the condition:
\begin{equation}
f^{\text{model}}(0,T) = f^M(0,T) \quad \forall\, T > 0
\end{equation}
The model-implied forward rate satisfies the Riccati-type equation:
\begin{equation}
\frac{\partial f^{\text{model}}(0,T)}{\partial T} = \theta(T) - a\,f^{\text{model}}(0,T) - \frac{\sigma^2}{2}\,\mathcal{B}(0,T)^2 \cdot \frac{\partial}{\partial T}\mathcal{B}(0,T)^{-1}
\end{equation}
Using $\mathcal{B}(0,T) = (1-e^{-aT})/a$, one computes $\frac{\partial}{\partial T}[\mathcal{B}(0,T)^2] = 2\mathcal{B}(0,T)\,e^{-aT}$, so:
\begin{equation}
\frac{\sigma^2}{2} \cdot 2\mathcal{B}(0,T)\,e^{-aT} = \frac{\sigma^2}{a}\left(1 - e^{-aT}\right)e^{-aT} = \frac{\sigma^2}{2a}\frac{\partial}{\partial T}\!\left(1-e^{-2aT}\right)
\end{equation}
Setting $f^{\text{model}}(0,T) = f^M(0,T)$ and solving for $\theta(T)$ then gives:
\begin{equation}
\boxed{\theta(t) = \frac{\partial f^M(0,t)}{\partial t} + a\, f^M(0,t) + \frac{\sigma^2}{2a}\left(1 - e^{-2at}\right)}
\end{equation}
The three terms have a clear interpretation: the first corrects for the slope of the forward curve, the second accounts for mean reversion pulling the rate back, and the third is a convexity adjustment from the stochastic volatility. As a sanity check: at $t=0$, we get $r(0) = f^M(0,0) = 2.03\%$ and $\theta(0) = 0.38\%$, consistent with a positively sloped yield curve.

% ===== SECTION 2 =====
\section{Multi-Measure Monte Carlo Pricing}

\subsection{Swaption Description}

We price a European payer swaption: exercise date $T = 2$ years, 10-year underlying swap with annual payments at $T_1 = 3, \ldots, T_{10} = 12$, strike $K = 2\%$, $\delta_i = 1$. The payoff at $T$ is:
\begin{equation}
V_T = A(T) \cdot \max\!\left(S(T) - K,\, 0\right)
\end{equation}
where $A(T) = \sum_{i=1}^{10} P(T, T+i)$ and $S(T) = (1 - P(T, T_{10})) / A(T)$.

\subsection{Hull-White Affine Bond Formula}

ZCB prices under Hull-White have the affine form $P(t,T) = \mathcal{A}(t,T)\,e^{-\mathcal{B}(t,T)\,r_t}$, with:
\begin{align}
\mathcal{B}(t,T) &= \frac{1-e^{-a(T-t)}}{a}, \qquad
\mathcal{A}(t,T) = \frac{P^M(0,T)}{P^M(0,t)} \exp\!\left(\mathcal{B}(t,T)\, f^M(0,t) - \frac{\sigma^2}{4a}(1-e^{-2at})\,\mathcal{B}(t,T)^2\right)
\end{align}

\subsection{Exact Simulation Scheme}

Under $\Q$ and $\Q^T$, the short-rate SDE has a known closed-form transition, so we can simulate each step exactly without any Euler approximation:
\begin{equation}
r_{t+\Delta t} = r_t\, e^{-a\Delta t} + \underbrace{\int_t^{t+\Delta t} \theta^\star(s)\, e^{-a(t+\Delta t - s)}\diff s}_{\mu_i \text{ (precomputed via \texttt{quad})}} + \sigma\sqrt{\frac{1-e^{-2a\Delta t}}{2a}} \cdot Z_i, \quad Z_i \sim \mathcal{N}(0,1)
\end{equation}
Under $\Q$, the discount factor $e^{-\int_0^T r_s\diff s}$ is approximated by the trapezoidal rule. This introduces a quadrature error of order $O(\Delta t^2)$, which is negligible compared to Monte Carlo noise at any reasonable step count.

\subsection{Risk-Neutral Measure $\Q$}

\textbf{Numeraire:} bank account $B_t$. \textbf{Drift:} $\theta^\star = \theta(t)$. \textbf{Pricing:} $V_0 = \E^{\Q}[e^{-\int_0^T r_s\diff s} V_T]$.

\begin{algorithm}[H]
\caption{Monte Carlo under $\Q$ (exact scheme)}
\begin{algorithmic}[1]
\State Precompute $\mu_i = \int_{t_i}^{t_{i+1}} \theta(s)\,e^{-a(t_{i+1}-s)}\diff s$ for all $i$
\State Initialize $r_0^{(k)} = f^M(0,0)$
\For{$i = 0, \ldots, n_{\text{steps}}-1$}
    \State $r_{i+1}^{(k)} \leftarrow r_i^{(k)}\,e^{-a\Delta t} + \mu_i + \text{vol}\cdot Z_i^{(k)}$; \quad $I^{(k)} \mathrel{+}= \tfrac{1}{2}(r_i^{(k)} + r_{i+1}^{(k)})\,\Delta t$
\EndFor
\State $\hat{V}_0 = \frac{1}{N}\sum_k e^{-I^{(k)}} \cdot A(r_T^{(k)})\cdot\max(S(r_T^{(k)})-K, 0)$
\end{algorithmic}
\end{algorithm}

\subsection{T-Forward Measure $\Q^T$}

\textbf{Numeraire:} ZCB $P(t,T)$. Under $\Q^T$, the numeraire is $P(t,T)$, so the Radon-Nikodym density is:
\begin{equation}
\frac{\diff\Q^T}{\diff\Q}\bigg|_t = \frac{P(t,T)/P(0,T)}{B_t} = \frac{P(t,T)}{P(0,T)\,B_t}
\end{equation}
In the Hull-White model, $P(t,T) = \mathcal{A}(t,T)\,e^{-\mathcal{B}(t,T)\,r_t}$, so by Itô's lemma:
\begin{equation}
\diff\ln P(t,T) = \left(\ldots\right)\diff t - \mathcal{B}(t,T)\,\sigma\,\diff W_t^{\Q}
\end{equation}
The Girsanov kernel is therefore $\lambda(t) = -\sigma\,\mathcal{B}(t,T)$, and the new Brownian motion under $\Q^T$ is:
\begin{equation}
\diff W_t^{\Q^T} = \diff W_t^{\Q} + \sigma\,\mathcal{B}(t,T)\,\diff t
\end{equation}
Substituting into the $\Q$-dynamics of $r_t$ gives:
\begin{align}
\diff r_t &= (\theta(t) - a\,r_t)\diff t + \sigma\,\diff W_t^{\Q} \\
          &= \bigl(\theta(t) - a\,r_t - \sigma^2\,\mathcal{B}(t,T)\bigr)\diff t + \sigma\,\diff W_t^{\Q^T}
\end{align}
The drift correction $-\sigma^2\,\mathcal{B}(t,T)$ is what we called $\theta^\star(s) - \theta(s)$. Because $\mathcal{B}(t,T)$ is deterministic, this term can be precomputed, which is why the exact simulation scheme still applies under $\Q^T$.

The main advantage of this measure change is practical: since $P(t,T)$ is the numeraire, the price formula simplifies to:
\begin{equation}
V_0 = P^M(0,T) \cdot \E^{\Q^T}[V_T]
\end{equation}
There is no stochastic discount factor to simulate — the randomness is entirely absorbed into the drift of $r_t$.

\begin{algorithm}[H]
\caption{Monte Carlo under $\Q^T$ (exact scheme)}
\begin{algorithmic}[1]
\State Precompute $\mu_i^T = \int_{t_i}^{t_{i+1}} (\theta(s) - \sigma^2 \mathcal{B}(s,T))\,e^{-a(t_{i+1}-s)}\diff s$
\State Simulate paths; \quad $\hat{V}_0 = P^M(0,T) \cdot \frac{1}{N}\sum_k V_T^{(k)}$
\end{algorithmic}
\end{algorithm}

\subsection{Annuity Measure $\Q^A$}

\textbf{Numeraire:} annuity $A_{n,m}(t) = \sum_j \tau_j P(t,T_j)$. Under this measure, the forward swap rate $S(t)$ is a martingale, which is natural for swaption pricing. The required drift correction is:
\begin{equation}
\theta^\star(s) = \theta(s) - \sigma^2 \sum_{j} \frac{\delta\, P(s,T_j)}{A(s,r_s)}\,\mathcal{B}(s,T_j)
\end{equation}
Since this term depends on the current path of $r_s$, it cannot be precomputed and requires an Euler step. The price is then $V_0 = A(0)\cdot\E^{\Q^A}[\max(S(T)-K,0)]$.

\begin{algorithm}[H]
\caption{Monte Carlo under $\Q^A$ (hybrid Euler scheme)}
\begin{algorithmic}[1]
\For{$i = 0, \ldots, n_{\text{steps}}-1$}
    \State Compute $P(t_i, T_j, r_i^{(k)})$ for all $j$, all $k$
    \State $c_i^{(k)} \leftarrow \sigma^2 \sum_j \frac{\delta\, P(t_i,T_j)}{A(t_i,r_i^{(k)})} \mathcal{B}(t_i,T_j)$
    \State $r_{i+1}^{(k)} \leftarrow r_i^{(k)} + (\theta(t_i) - a\,r_i^{(k)} - c_i^{(k)})\,\Delta t + \sigma\sqrt{\Delta t}\, Z_i^{(k)}$
\EndFor
\State $\hat{V}_0 = A(0)\cdot\frac{1}{N}\sum_k \max(S_T^{(k)} - K, 0)$
\end{algorithmic}
\end{algorithm}

\subsection{Monte Carlo comparison across numeraires}

We price the same payer swaption under the three measures $Q$, $Q^T$ and $Q^A$.
In theory, the arbitrage-free price should be invariant to the choice of numeraire.
Our Monte Carlo results confirm this property:

\begin{table}[H]
\centering
\begin{tabular}{lccc}
\toprule
Measure & Price & Std. Error & Variance of estimator \\
\midrule
$Q$   & 0.099006 & 0.000158 & $2.4964\times 10^{-8}$ \\
$Q^T$ & 0.099004 & 0.000161 & $2.5921\times 10^{-8}$ \\
$Q^A$ & 0.099001 & 0.000174 & $3.0276\times 10^{-8}$ \\
\bottomrule
\end{tabular}
\caption{Monte Carlo prices under the three numeraires. The variance of the estimator is computed as $(\text{Std. Error})^2$.}
\end{table}

The three prices are almost identical, which is consistent with change-of-numeraire theory:
changing the numeraire does not change the theoretical price in an arbitrage-free framework.
The differences therefore come from Monte Carlo efficiency rather than from valuation inconsistency.

The main difference appears in the estimator variance. In our implementation, the risk-neutral
measure $Q$ produces the smallest standard error, followed by $Q^T$, while $Q^A$ gives the
largest one. At first sight, this may seem surprising, since under the annuity measure the
swaption payoff takes the clean form
\[
V_0 = A(0)\,\mathbb{E}^{Q^A}\big[(S(T)-K)^+\big].
\]
However, the observed variance ranking is not determined by the numeraire alone.

A first reason is that under $Q$, the simulated quantity contains both the random discount factor
and the random annuity term. These terms may partially offset the variability of the option payoff:
when rates increase, the swap rate and the payoff tend to increase, while discount factors and bond
prices tend to decrease. This compensation effect can reduce the overall dispersion of the Monte Carlo
integrand.

A second and more important reason comes from the numerical implementation. Under $Q$ and $Q^T$,
the short-rate dynamics remain Gaussian affine and admit an exact discretization. Under $Q^A$,
by contrast, the drift correction depends on the current annuity weights, which are themselves
functions of zero-coupon bond prices and therefore of the current short rate. The process is no
longer affine Gaussian, so an Euler discretization must be used. As a consequence, the empirical
standard error under $Q^A$ reflects both the numeraire choice and an additional discretization effect.

Therefore, the results should be interpreted as follows: the equality of prices validates the
change-of-numeraire framework, while the variance ranking mainly shows that, in our implementation,
the risk-neutral measure is numerically the most efficient. This does not mean that $Q$ is
theoretically always the best numeraire, but rather that it is the most stable one in the present
short-rate simulation framework.
% ===== SECTION 3 =====
\section{Numerical Implementation: Trinomial Trees}

\subsection{Construction of the Hull-White Trinomial Tree}

Following Hull \& White (1993), we build a recombining trinomial tree on the centered process $x_t = r_t - \alpha(t)$, which satisfies $\diff x_t = -a\,x_t\,\diff t + \sigma\,\diff W_t$. The node spacing is:
\begin{equation}
\Delta x = \sigma\sqrt{3\,\Delta t}
\end{equation}
This choice minimizes discretization error. Beyond $j_{\max} = \lceil 0.184 / (a\,\Delta t) \rceil$ nodes from the center, the branching geometry is modified to keep all probabilities positive.

\subsubsection{Transition probabilities}

At a standard interior node $j$, setting $\eta = a\,j\,\Delta t$:
\begin{equation}
p_u = \frac{1}{6} + \frac{\eta^2 + \eta}{2}, \qquad p_m = \frac{2}{3} - \eta^2, \qquad p_d = \frac{1}{6} + \frac{\eta^2 - \eta}{2}
\end{equation}
At boundary nodes $j = \pm j_{\max}$, the branching direction shifts by $\mp 1$ to keep the tree within bounds while preserving the first two conditional moments.

\subsubsection{Node-shifting $\alpha(t)$ to fit the yield curve}

The shift $\alpha(t_i)$ is computed recursively using the Arrow-Debreu prices $Q(i,j)$ — the time-0 value of a security that pays 1 at node $(i,j)$ and 0 elsewhere. The no-arbitrage condition gives:
\begin{equation}
\alpha(t_i) = \frac{1}{\Delta t}\left[\ln\!\sum_j Q(i,j)\,e^{-j\,\Delta x\,\Delta t} - \ln P^M(0, t_{i+1})\right]
\end{equation}
With this choice, $r(t_i, j) = \alpha(t_i) + j\,\Delta x$ exactly matches the market ZCB prices at each step.

\subsubsection{Yield curve fit verification}

\begin{table}[H]
\centering
\begin{tabular}{cccc}
\toprule
\textbf{Maturity} & \textbf{$P^{\text{NSS}}(0,T)$} & \textbf{$P^{\text{tree}}(0,T)$} & \textbf{Error (bps)} \\
\midrule
0.5  & 0.98978 & 0.98978 & $<0.001$ \\
1.0  & 0.97909 & 0.97909 & $<0.001$ \\
2.0  & 0.95757 & 0.95757 & $<0.001$ \\
5.0  & 0.89111 & 0.89111 & $<0.001$ \\
10.0 & 0.75390 & 0.75390 & $<0.001$ \\
\bottomrule
\end{tabular}
\caption{Tree vs.\ NSS market ZCB prices ($N = 400$)}
\end{table}

The tree fits the NSS curve to machine precision. A few representative maturities are shown; the full grid was verified numerically.

\subsection{Swaption Pricing by Backward Induction}

At the exercise date $T_{\text{ex}}$, each node $j$ has a known short rate $r(T_{\text{ex}}, j) = \alpha(T_{\text{ex}}) + j\,\Delta x$. We compute the swaption payoff at each node using the Hull-White ZCB formula:
\begin{equation}
V(T_{\text{ex}}, j) = A(T_{\text{ex}}, j)\cdot\max\!\left(S(T_{\text{ex}}, j) - K,\, 0\right)
\end{equation}
We then roll back from $i = i_{\text{ex}}$ to $i = 0$ using:
\begin{equation}
V(t_i, j) = e^{-r(t_i,j)\,\Delta t}\!\left[p_u V(t_{i+1}, j_u) + p_m V(t_{i+1}, j_m) + p_d V(t_{i+1}, j_d)\right]
\end{equation}

\begin{algorithm}[H]
\caption{Trinomial Tree Pricing}
\begin{algorithmic}[1]
\State Build tree: compute $\alpha(t_i)$ and Arrow-Debreu prices $Q(i,j)$ for $i = 0,\ldots,N$
\State At $i = i_{\text{ex}}$: compute $V(T_{\text{ex}}, j) = A_j \cdot \max(S_j - K, 0)$ using $P_{\text{HW}}$
\For{$i = i_{\text{ex}}-1$ down to $0$}
    \For{each node $j$ at step $i$}
        \State $V(t_i,j) \leftarrow e^{-r(t_i,j)\Delta t}(p_u V_{j_u} + p_m V_{j_m} + p_d V_{j_d})$
    \EndFor
\EndFor
\State \Return $V(0, 0)$
\end{algorithmic}
\end{algorithm}

\subsection{Comparison with Monte Carlo Results}

\begin{table}[H]
\centering
\begin{tabular}{lccc}
\toprule
\textbf{Method} & \textbf{Price} & \textbf{$\pm 2\sigma$} & \textbf{Rel.\ diff vs.\ tree} \\
\midrule
Trinomial Tree ($N=400$) & 0.098969 & N/A      & --- \\
Risk-Neutral MC ($\Q$)   & 0.099006 & 0.000316 & $+0.037\%$ \\
Forward MC ($\Q^T$)      & 0.099004 & 0.000323 & $+0.035\%$ \\
Annuity MC ($\Q^A$)      & 0.099001 & 0.000348 & $+0.033\%$ \\
\bottomrule
\end{tabular}
\caption{Swaption price: trinomial tree vs.\ Monte Carlo (50,000 paths, 400 steps)}
\end{table}

All four results are within $0.04\%$ of each other. The small gap between the tree and the MC prices is well inside the MC confidence intervals, so there is no evidence of a systematic bias in either method.

\subsection{Convergence Analysis}

\subsubsection{Trinomial tree convergence}

\begin{table}[H]
\centering
\begin{tabular}{rcc}
\toprule
\textbf{N} & \textbf{Price} & \textbf{$\Delta$ vs $N=400$} \\
\midrule
25  & 0.100330 & $+0.001361$ \\
50  & 0.099337 & $+0.000368$ \\
75  & 0.099021 & $+0.000052$ \\
100 & 0.099536 & $+0.000567$ \\
150 & 0.099133 & $+0.000164$ \\
200 & 0.098932 & $-0.000037$ \\
300 & 0.098959 & $-0.000010$ \\
400 & 0.098969 & ---          \\
\bottomrule
\end{tabular}
\caption{Trinomial tree convergence with respect to number of time steps $N$}
\end{table}

The price stabilizes around $N=200$--$300$, with an error that becomes small as the grid is refined. The non-monotone jump at $N=100$ is typical of trinomial trees: it happens when the exercise date falls between grid points, which changes the local geometry of the backward induction.

\subsubsection{Monte Carlo convergence}

\begin{table}[H]
\centering
\begin{tabular}{rcccc}
\toprule
\textbf{$N_{\text{steps}}$} & \textbf{$\Delta t$} & \textbf{Price (MC)} & \textbf{StdErr} \\
\midrule
 25  & 0.0800 & 0.098685 & 0.000158 \\
 50  & 0.0400 & 0.098836 & 0.000158 \\
100  & 0.0200 & 0.098894 & 0.000158 \\
200  & 0.0100 & 0.098914 & 0.000157 \\
400  & 0.0050 & 0.099006 & 0.000158 \\
600  & 0.0033 & 0.098897 & 0.000158 \\
\bottomrule
\end{tabular}
\caption{Monte Carlo convergence under $\Q$ (50,000 paths) with respect to number of steps}
\end{table}

The MC price barely moves as we refine the time grid — the standard error is essentially constant, and all prices fall within each other's confidence intervals. This is consistent with the exact transition scheme: there is no SDE discretization to refine away, so finer grids only add computation time without improving accuracy.

\subsubsection{Comparison of convergence behavior}

The two methods converge very differently. The MC price is stable from $N_{\text{steps}} = 25$ onward because the short-rate transition is simulated exactly — there is simply no discretization error to reduce. Residual fluctuations are pure Monte Carlo noise of order $O(1/\sqrt{N_{\text{paths}}})$. The tree price, by contrast, oscillates at coarse grid sizes, but the oscillation amplitude decreases as the grid is refined, before settling around $N \approx 200$. In terms of speed, the tree is much faster: it gives a stable result in under 0.5 seconds, while 50,000 MC paths take several seconds.

\subsection{Sensitivity to the Mean-Reversion Parameter $a$}

\begin{table}[H]
\centering
\begin{tabular}{rcc}
\toprule
\textbf{$a$} & \textbf{Price} & \textbf{$\Delta$ vs calibrated} \\
\midrule
0.01  & 0.098995 & $+0.000026$ \\
0.05  & 0.098975 & $+0.000006$ \\
0.10  & 0.098963 & $-0.000006$ \\
0.15  & 0.098975 & $+0.000006$ \\
0.20  & 0.099018 & $+0.000049$ \\
0.30  & 0.099225 & $+0.000256$ \\
0.40  & 0.099288 & $+0.000319$ \\
0.50  & 0.098401 & $-0.000568$ \\
0.60  & 0.097644 & $-0.001325$ \\
0.70  & 0.098098 & $-0.000871$ \\
\bottomrule
\end{tabular}
\caption{Swaption price sensitivity to $a$ (trinomial tree, $N=400$, $\sigma$ fixed)}
\end{table}

For $a \lesssim 0.40$, the price barely moves — deviations stay below 0.3\%. There is a local maximum around $a \approx 0.35$--$0.40$, after which the price starts to fall.

\textbf{Interpretation.} A higher $a$ means the short rate reverts faster toward the path imposed by $\theta(t)$. This compresses the distribution of $r_T$, which tends to reduce the probability of the swaption being deep in the money. At the same time, $\theta(t)$ re-adjusts to keep the yield curve fit, partially cancelling that effect. Near the calibrated value $a = 0.1367$, the price is quite flat, which is reassuring: small errors in estimating $a$ do not significantly affect the result. For $a \gtrsim 0.45$ the price drops more sharply, but this region should be interpreted with caution — part of the behavior is likely numerical, since large $a$ values reduce $j_{\max} = \lceil 0.184/(a\,\Delta t)\rceil$ and make the tree coarser.

% ===== SECTION 4 =====
\section{Critical Analysis}

\subsection{The Assumption $P = \Q$ on $(a, \sigma)$}

We assumed throughout that $(a, \sigma)$ are the same under the historical measure $P$ and the risk-neutral measure $\Q$. This is a convenient simplification — it lets us estimate the parameters from historical data without needing market option prices — but it is not necessarily true. In practice, interest rate risk premia create a wedge between $P$ and $\Q$. Under $P$:
\begin{equation}
\diff r_t = (\theta^P(t) - a\,r_t)\diff t + \sigma\,\diff W_t^P
\end{equation}
Girsanov's theorem gives $\diff W_t^P = \diff W_t^{\Q} + \lambda(t)\diff t$. By fixing $a$ and $\sigma$ across measures, we are absorbing all of the measure change into $\theta(t)$ and ignoring the possibility that the mean-reversion speed itself differs between the two worlds. This can distort long-maturity option prices when the market prices a significant volatility risk premium.

\subsection{Historical vs.\ Market-Implied Volatility}

The $\sigma = 0.0067$ we estimated is a \textit{historical volatility} — it reflects how much the \texteuro STR actually moved over the past six years. \textit{Implied volatility} is something different: it is the value of $\sigma$ that makes the Hull-White formula match an observed market swaption price (via the Jamshidian decomposition). The two need not agree. Implied volatility incorporates market expectations and a risk premium; historical volatility does not. In stress periods, implied volatility can be significantly higher, which means a model calibrated on historical data will tend to underprice options. Our price of $0.0990$ should therefore be read as a lower bound relative to what the market might quote.

\subsection{Calibration to Market Swaption Prices}

If market swaption prices were available, we would calibrate very differently:
\begin{enumerate}
    \item \textbf{Input:} a grid of implied Black volatilities across exercise maturities and swap tenors (the ``vol cube'').
    \item \textbf{Target:} the Hull-White analytical swaption price (Jamshidian decomposition) as a function of $(a, \sigma)$.
    \item \textbf{Fit:} minimize the squared pricing error across the swaption grid.
    \item \textbf{Output:} risk-neutral $(a, \sigma)$ that best match observed market prices.
\end{enumerate}
This gives parameters that are directly consistent with market option prices, which is what matters for hedging and pricing derivatives. The trade-off is that we lose the connection to historical dynamics. A natural extension is to allow $\sigma = \sigma(t)$, which adds enough degrees of freedom to fit the volatility surface across maturities — this is the standard approach in practice.

\end{document}